\section{Discussion}
\label{sec:disc}

For benchmark authors, SRE teams, and method developers, the fitted
chain provides a shared reliability object at the trace level: it
supports first-passage metrics with uncertainty, horizon and
perturbation queries, and diagnostics that can reject an inadequate
abstraction. SS9 (Table~\ref{tab:heldout-ss9}) shows that, once the
trace-to-chain step is specified and checked, $\hat M$ tracks held-out
first-passage behavior to within $L_\infty^{\mathrm{RDC}}\le 0.053$
across seven frameworks. This does not replace benchmarks or incident
analysis; it shows how benchmark-style trace corpora can support
auditable reliability answers while also flagging when the abstraction
is too weak for them. The practical requirement is to fit the chain,
report uncertainty, run GoF, and withhold first-passage conclusions
when the checks fail.

The real-data studies make this concrete and show the audit is
\emph{selective} rather than all-or-nothing: across SWE-bench (SS10,
SS14) and $\tau$-bench (SS15) it accepts the asymptotic reliability but
rejects the first-order RDC timing, while SS11--SS13
(\S\ref{sec:robustness}) show the verdict survives the clustering
choice even where the point estimate does not. The value of the fitted
chain on real traces is thus real but \emph{partial}, and the audit is
what makes it safe by saying which part to trust.

The clearest next step, identified by the data rather than assumed, is a
duration-aware model: SS14 localizes the misfit to the success-time
distribution, so semi-Markov holding times (or a learned phase-type
duration law) are the indicated route to restoring horizon queries.
Further directions include HMM reliability under partial observability,
online estimation of $Q$, integration with PRISM and Storm, and
mixture-chain models for cross-trial correlation~\cite{spaeh2024}.

%% ================================================================
